% Chapter 2 - State of the art and technologies.
% Sources: docs/technologies.md + new text with IEEE bibliography.
% Target length: 8-10 pages.

\chapter{State of the art and technologies}
\label{ch:estado-arte}

% TODO: chapter intro, ~1 page.
% "This chapter places UnivOrch in the landscape of existing
% orchestration solutions. We first review the classical categories
% (hypervisor management, IaC, container orchestration), then narrow
% down what the teaching case requires and why off-the-shelf products
% do not fully fit. We close with the concrete technologies adopted
% for the implementation."

\section{Virtualisation and orchestration}
\label{sec:virtualizacion-orquestacion}

% TODO: draft. Ideas:
% - Conceptual difference between "hypervisor", "hypervisor manager"
%   (vCenter) and "orchestration" (vRA, OpenStack, Kubernetes).
% - Declarative vs imperative model: Terraform as a reference.
% - Conceptual parallel with Kubernetes: control plane + declarative
%   API + reconciliation. UnivOrch borrows the resource model without
%   inheriting Kubernetes' complexity.

\section{Existing solutions}
\label{sec:soluciones-existentes}

% TODO: draft. Comparison table at the end.
% Products to discuss:
% - VMware vSphere (vCenter): the de-facto standard on VMware.
%   Paid, proprietary.
% - VMware vRealize Automation (vRA): VMware's orchestration layer.
%   Complex, expensive, proprietary.
% - OpenStack: open source IaaS orchestrator. Excessively complex for
%   small teaching environments.
% - Proxmox VE: management UI for Proxmox. Bound to Proxmox.
% - Foreman: open source provisioning. More oriented to bare metal.
% - Cloud-init / Ansible: configuration, not orchestration.
%   Complementary, not substitutes.
% - Kubernetes: container orchestration, not VM orchestration.
%   Reference model.
%
% For each: one paragraph stating what it solves and what does not
% match the teaching case.

\section{The teaching case: specific needs}
\label{sec:caso-docente}

% TODO: draft. Ideas:
% - One course == a pyramid of N copies of a template.
% - Recycling at the end of the term: bulk undeploy.
% - Isolation between students (network, IPs, snapshots).
% - The lecturer should not need to be a VMware or networking expert.
% - Heterogeneity: VMware in the classroom lab, Proxmox in the
%   research lab. Same conceptual model expected.
% - This leaves a gap: a hypervisor-agnostic orchestrator with
%   vocabulary fitting the teaching environment.

\section{Starting point: supervisor's previous work}
\label{sec:punto-partida}

% TODO: draft. Ideas:
% - esxobjects: Python library wrapping the vSphere SOAP API into a
%   friendlier interface. Bound to VMware.
% - yamlinfr: reads YAML descriptions and deploys to vSphere using
%   esxobjects. Embryonic declarative model.
% - UnivOrch generalises the model to any hypervisor, rewrites the
%   connectors from scratch under TDD, and adds a REST layer as a
%   public boundary.
% - The Development chapter includes a qualitative comparison with
%   both libraries.

\section{Selected technologies and rationale}
\label{sec:tecnologias-seleccionadas}

% TODO: summary table at the top, then justify each block.
% Source: docs/technologies.md (full content).
% Blocks:
% - Language: Python 3.12.
% - Web framework: FastAPI + uvicorn.
% - HTTP client: httpx (with injection for tests).
% - Storage: TinyDB (PoC) -> MongoDB (future). Document database.
% - Modelling: Pydantic v2.
% - CLI: cmd2.
% - Quality: ruff, mypy strict, pytest + Hypothesis.
% - Packaging: multi-stage Docker, ghcr.io.
% - Hypervisor libraries (future): pyvmomi, proxmoxer.
%
% For each: what it is used for, why chosen over alternatives, what
% was discarded and why.
